\documentclass[10pt]{beamer}

% Theme and Color setup
\usetheme{Madrid}
\usecolortheme{beaver}
\setbeamertemplate{navigation symbols}{}
\setbeamertemplate{caption}[numbered]

% Force Theorem blocks to be red
\definecolor{darkred}{RGB}{160, 0, 0}
\setbeamercolor{block title}{bg=darkred, fg=white}
\setbeamercolor{block body}{bg=red!10, fg=black}

% Packages
\usepackage{amsmath, amssymb, amsthm}
\usepackage{mathrsfs}
\usepackage{mathtools}
\usepackage{tikz}
\usepackage{lmodern}
\usefonttheme{professionalfonts}
% \usepackage{fontspec}
% \setmainfont{TeX Gyre Heros} % or another modern font
\usetikzlibrary{arrows.meta, positioning}

\newtheorem{intuition}{Intuition}
\newtheorem{motivation}{Motivation}

\title[Fully Homomorphic Encryption over the Integers]{Fully Homomorphic Encryption over the Integers}
\subtitle{Paper by van Dijk, Gentry, Halevi, and Vaikuntanathan (2010)}
\author[Nirav, Kunal]{Nirav Bhattad (23B3307), Kunal Chaudhari (22B1060)}
\date{\today}

\begin{document}

\begin{frame}
    \titlepage
\end{frame}

\begin{frame}{Outline}
    \tableofcontents[hideallsubsections]
\end{frame}

\section{The FHE Dream and The Bootstrapping Blueprint}

\subsection{Motivation \& The Integer Paradigm}

\begin{frame}{Motivation: The Cloud Computing Conundrum}
    \begin{motivation}
    We possess robust cryptographic solutions for data \textbf{in transit} (TLS/SSL) and data \textbf{at rest} (AES). However, computing on that data traditionally requires decryption, creating a massive vulnerability when outsourcing to untrusted cloud servers.
    \end{motivation}
    \pause
    
    \vspace{0.3cm}
    \textbf{The Fully Homomorphic Encryption (FHE) Goal:}
    Allow an untrusted party to take encrypted inputs, perform arbitrary, unbounded computations on them, and return an encrypted result. 
    \pause
    
    \vspace{0.3cm}
    When decrypted, this result must match exactly what the data owner would have obtained by computing on the plaintext locally.
\end{frame}

\begin{frame}{The Quest for the Simplest Secure Scheme}
    \textit{``What is the simplest encryption scheme for which one can hope to achieve security?"}
    \pause
    
    \vspace{0.3cm}
    \textbf{The Caesar Cipher Analogy:}
    \begin{itemize}
        \item Shift `A' by 3 and `B' by 3. You can trivially add those shifts together.
        \item It is perfectly (and ironically) homomorphic under addition. 
        \item \textbf{Fatal Flaw:} It is completely insecure, because the ciphertexts are just shifted versions of the plaintexts, making them easily breakable by frequency analysis.
    \end{itemize}
    \pause

    \vspace{0.3cm}
    \textbf{The Paradigm Shift (2009 vs. 2010):}
    \begin{itemize}
        \item Craig Gentry's breakthrough 2009 FHE scheme utilized ideal lattices over polynomial rings -- using conceptually and mathematically heavy machinery like Algebraic Number Theory and Algebraic Geometry.
        \pause
        \item \textbf{This Paper:} Asks if we can achieve the exact same FHE blueprint using only elementary addition and multiplication over the integers. 
    \end{itemize}
\end{frame}

\subsection{Formalizing Fully Homomorphic Encryption}

\begin{frame}{Formalizing the FHE Dream}
    To prove a scheme works, we must rigorously define it. A homomorphic public key encryption scheme $\mathcal{E}$ consists of four Probabilistic Polynomial-Time (PPT) algorithms:
    
    \pause
    $$\mathcal{E} = (\text{KeyGen}, \text{Encrypt}, \text{Decrypt}, \textbf{Evaluate})$$
    
    \vspace{0.3cm}
    \textbf{The Evaluate Algorithm:}
    Takes as input a public key $pk$, a boolean circuit $C$, and a tuple of ciphertexts $\vec{c}=\langle c_{1},...,c_{t}\rangle$ (corresponding to the $t$ input bits of $C$), and outputs a new evaluated ciphertext $c_{\text{out}}$.
\end{frame}

\begin{frame}{Correctness and Homomorphism}
    \begin{definition}[Correct Homomorphic Decryption]
    The scheme $\mathcal{E}$ is correct for a given $t$-input circuit $C$ if, for any key-pair $(sk,pk)\leftarrow \text{KeyGen}(\lambda)$, any plaintexts $m_{1},...,m_{t}\in\{0,1\}$, and any ciphertexts $c_{i}\leftarrow \text{Encrypt}(pk,m_{i})$, the following holds with overwhelming probability:
    $$\text{Decrypt}(sk,\text{Evaluate}(pk,C,\vec{c}))=C(m_{1},...,m_{t})$$
    \end{definition}
    \pause

    \vspace{0.3cm}
    \begin{definition}[Fully Homomorphic Encryption]
    The scheme $\mathcal{E}$ is homomorphic for a class $\mathcal{C}$ of circuits if it is correct for all circuits $C\in\mathcal{C}$. The scheme $\mathcal{E}$ is \textbf{fully homomorphic} if it is correct for \textit{all} boolean circuits (i.e., circuits of arbitrary depth and size).
    \end{definition}
\end{frame}

\begin{frame}{The ``Trivial" Solution and Circuit Privacy}
    \textbf{Loophole:} Can we trivially satisfy the definition of FHE right now? 
    \pause
    
    \textit{Yes.} An untrusted server could just take the circuit $C$, concatenate it to the raw input ciphertexts, and send it back. The user decrypts the inputs, reads $C$, and computes it locally.
    \pause
    
    % \vspace{0.3cm}
    To ban this useless scheme, cryptography enforces two strict properties. The first is \textbf{Circuit Privacy}.
    \pause

    \begin{definition}[Circuit Privacy]
    The ciphertext output by the $\text{Evaluate}$ algorithm must not reveal anything about the circuit $C$ that was evaluated, beyond the final output value itself -- even to the data owner who holds the secret key.
    \end{definition}
    \pause

    % \begin{intuition}
    % If the server just attaches $C$ to the ciphertext, privacy is completely violated. In our integer scheme, even the bit-length of the output ciphertext could leak how many multiplication gates were evaluated. To achieve circuit privacy, we must "wipe the history" of the computation by re-randomizing the final ciphertext -- specifically by adding a fresh ``high noise" encryption of zero before sending it back.
    % \end{intuition}
\end{frame}

\begin{frame}{The Compactness Constraint}
    The second (and far more difficult) property that rules out the trivial solution is \textbf{Compactness}.
    \pause
    
    \vspace{0.3cm}
    \textbf{The Goal:} We must physically prevent the server from passing the computational burden back to the user. The server must be forced to actually do the math and compress the answer.
    \pause

    \begin{definition}[Compactness]
    A homomorphic scheme $\mathcal{E}$ is compact if there exists a fixed polynomial bound $b(\lambda)$ such that the size of the ciphertext output by $\text{Evaluate}(pk, C, \vec{c})$ is strictly bounded by $b(\lambda)$ bits, \textbf{independently of the size or depth of $C$}.
    \end{definition}
    \pause

    \begin{intuition}[The Root of the FHE Challenge]
    Compactness is the true bottleneck of FHE. We cannot let the ciphertext grow proportionally to the circuit depth. We must compress the output so the user experiences $\mathcal{O}(1)$ decryption time, regardless of whether the cloud server evaluated a 10-gate circuit or a 10-billion-gate circuit. This strict bounding is exactly what forces us to manage noise, and ultimately, what forces us to use Bootstrapping.
    \end{intuition}
\end{frame}

\subsection{The Bootstrapping Blueprint}

\begin{frame}{The Crisis of Noise}
    \begin{intuition}[The Noise Bottleneck]
    In all known FHE schemes, ciphertexts contain ``noise" to ensure security. 
    \begin{itemize}
        \item Homomorphic addition adds the noise.
        \item Homomorphic multiplication \textbf{multiplies} the noise.
    \end{itemize}
    Because multiplication squares the noise, evaluating a deep circuit causes the noise to explode past the decryption boundary, permanently destroying the plaintext.
    \end{intuition}
    \pause
    
    \vspace{0.3cm}
    This limits us to a \textit{Somewhat} Homomorphic Encryption (SHE) scheme. How do we compute infinitely if noise acts as a strict depth limit?
\end{frame}

\begin{frame}{Defining Bootstrappability}
    The answer is Gentry's Bootstrapping Blueprint. First, we define a specific class of circuits.
    \pause

    \vspace{0.3cm}
    \begin{definition}[Augmented Decryption Circuits]
    Let $\mathcal{E}$ be an encryption scheme where decryption is implemented by a boolean circuit. The set of augmented decryption circuits, $D_{\mathcal{E}}$, consists of two circuits: 
    \begin{enumerate}
        \item Decrypt two ciphertexts and ADD the results.
        \item Decrypt two ciphertexts and MULTIPLY the results.
    \end{enumerate}
    \end{definition}
    \pause

    \vspace{0.3cm}
    \begin{definition}[Bootstrappable Encryption]
    Let $\mathcal{C}_{\mathcal{E}}$ be the set of permitted circuits a scheme can evaluate before noise exceeds the bounds. A scheme is \textbf{bootstrappable} if it can evaluate its own augmented decryption circuits:
    $$D_{\mathcal{E}} \subseteq \mathcal{C}_{\mathcal{E}}$$
    \end{definition}
\end{frame}

\begin{frame}{Gentry's Bootstrapping Theorem}
    \begin{theorem}[Gentry's Transformation]
    There is an efficient transformation that, given a description of a bootstrappable scheme $\mathcal{E}$, outputs a description of a \textbf{Fully Homomorphic} encryption scheme capable of evaluating circuits of any depth.
    \end{theorem}
    \pause
    
    \vspace{0.3cm}
    \begin{intuition}[How Bootstrapping ``Refreshes" Ciphertexts]
    \begin{enumerate}
        \item Take a ciphertext $c$ whose noise is dangerously high.
        \item Encrypt $c$ \textit{again} under a new public key.
        \item Provide the server with an \textit{encrypted} secret key.
        \item The server homomorphically evaluates the decryption algorithm on the inner ciphertext.
    \end{enumerate}
    \pause
    \textbf{Result:} Because the decryption circuit is evaluated homomorphically, the output is a fresh ciphertext encrypting the exact same message, but with a fixed, smaller noise level determined solely by the depth of the decryption circuit itself!
    \end{intuition}
\end{frame}

\section{The Integer Engine and The Noise Crisis}
\subsection{The Symmetric \& Public-Key SHE Schemes}

\begin{frame}{Building the Engine: The Symmetric ``Toy" Scheme}
    Before we look at the complex public-key version, we start with a beautifully simple symmetric-key foundation.
    \pause
    
    \vspace{0.3cm}
    \textbf{Parameters:} Let $\eta$ be the bit-length of the secret key, and $\rho$ be the bit-length of the noise.
    \pause
    
    \begin{block}{The Symmetric Algorithms}
    \begin{itemize}
        \item \textbf{KeyGen:} The secret key $p$ is a randomly chosen odd integer $p \in [2^{\eta-1}, 2^{\eta})$.
        \item \textbf{Encrypt($p, m$):} To encrypt a bit $m \in \{0,1\}$, choose a large random integer $q$ and a small random noise integer $r$ (where $|r| < 2^{\rho}$). 
        $$c = p \cdot q + 2r + m$$
        \item \textbf{Decrypt($p, c$):} Output $m' = (c \pmod p) \pmod 2$.
    \end{itemize}
    \end{block}
\end{frame}

\begin{frame}{Intuition: Why Symmetric Decryption Works}
    Let's break down the decryption equation: $m' = (c \pmod p) \pmod 2$.
    \pause
    
    \vspace{0.3cm}
    \begin{enumerate}
        \item Substitute $c$: $(p \cdot q + 2r + m \pmod p)$
        \pause
        \item Because $p \cdot q$ is an exact multiple of $p$, it vanishes perfectly. We are left with $(2r + m \pmod p)$.
        \pause
        \item \textbf{The Critical Condition:} The modulo $p$ operation only leaves $2r+m$ unchanged if its absolute value is strictly less than half the modulus. We strictly require:
        $$|2r + m| < p/2$$
        \pause
        \item Assuming the noise is small enough, we take $(2r + m) \pmod 2$. The $2r$ is even, so it vanishes, leaving exactly $m$.
    \end{enumerate}
    \pause
    
    \begin{block}{Motivation}
    To anyone without $p$, the ciphertext $c$ just looks like a massive, random integer. But algebraically, it is a point ``tethered" to a multiple of $p$ by a short, noisy string ($2r+m$).
    \end{block}
\end{frame}

\begin{frame}{Symmetric Decryption: A Visualizable Example}
    Let's put numbers to the theory.
    
    \vspace{0.2cm}
    \textbf{Setup:}
    \begin{itemize}
        \item Secret Key $p = 31$
        \item Message bit $m = 1$
    \end{itemize}
    \pause
    
    \textbf{Encrypt:}
    Choose a random $q = 5$ and a small noise $r = 2$.
    $$c = (31 \cdot 5) + 2(2) + 1 = 155 + 4 + 1 = 160$$
    The ciphertext is $160$.
    \pause
    
    \vspace{0.2cm}
    \textbf{Decrypt:}
    $$160 \pmod{31} = 5$$
    $$5 \pmod 2 = 1$$
    We successfully recovered $m = 1$. The noise ($2r+m = 5$) was safely smaller than $p/2$ ($15.5$).
\end{frame}

\begin{frame}{Transitioning to Public-Key}
    A true FHE scheme requires public-key infrastructure. We achieve this by publishing a massive list of ``encryptions of zero."
    \pause

    \vspace{0.2cm}
    \textbf{New Parameters (based on security parameter $\lambda$):}
    \begin{itemize}
        \item $\gamma$: Bit-length of the integers in the public key.
        \item $\tau$: Number of integers in the public key.
        \item $\rho'$: Secondary noise parameter for encryption $\rho' = \rho + \omega(\log \lambda)$.
    \end{itemize}
    \pause

    \vspace{0.2cm}
    \begin{block}{The Noise Distribution $\mathcal{D}_{\gamma,\rho}(p)$}
    To generate an ``encryption of zero", choose $q \leftarrow \mathbb{Z} \cap [0, 2^\gamma/p)$ and $r \leftarrow \mathbb{Z} \cap (-2^\rho, 2^\rho)$. 
    Output $x = pq + r$.
    \end{block}
\end{frame}

\begin{frame}{The Public-Key Algorithms}
    \textbf{KeyGen:}
    \begin{itemize}
        \item Sample $\tau + 1$ integers $x_0, x_1, \dots, x_\tau$ from $\mathcal{D}_{\gamma,\rho}(p)$. 
        \item Relabel so $x_0$ is the largest. Restart unless $x_0$ is odd and its noise $r_p(x_0)$ is even.
        \item Public key $pk = \langle x_0, x_1, \dots, x_\tau \rangle$.
    \end{itemize}
    \pause
    
    \vspace{0.3cm}
    \textbf{Encrypt($pk, m$):}
    Choose a random subset $S \subseteq \{1, 2, \dots, \tau\}$ and a noise integer $r \in (-2^{\rho'}, 2^{\rho'})$.
    $$c = \left[ m + 2r + 2\sum_{i \in S} x_i \right]_{x_0}$$
    \pause
    
    \begin{block}{Intuition}
    Every $x_i$ is effectively $0$ (plus noise). Summing a random subset of $x_i$ computes $0+0+0\dots$, creating a fresh, randomized encryption of zero. We multiply by 2 to keep the noise even, add our own fresh noise $2r$, and add $m$. The final modulo $x_0$ prevents the physical bit-size from growing infinitely large.
    \end{block}
\end{frame}

\subsection{Proof of Correctness}

\begin{frame}{The Structure of a Fresh Ciphertext}
    To prove this scheme is homomorphic, we first mathematically prove that a fresh public-key ciphertext has the exact same structure as our symmetric toy scheme.
    \pause
    
    \begin{lemma}
    Let $c \leftarrow \text{Encrypt}(pk, m)$. Then, $c = a\cdot p + (2b + m)$ for some integers $a$ and $b$ where the noise is bounded by $|2b+m| \le \tau 2^{\rho+2}$.
    \end{lemma}
    \pause
    
    \vspace{0.3cm}
    \textbf{Proof Sketch:}
    \begin{itemize}
        \item By definition, $c = m + 2r + 2\sum_{i \in S} x_i + k \cdot x_0$.
        \item We know every $x_i = q_i p + r_i$. Substitute this in:
        $$c = p \cdot \left(kq_0 + 2\sum_{i \in S} q_i\right) + \left(m + 2r + 2k \cdot r_0 + 2\sum_{i \in S} r_i\right)$$
        \item The first term is a multiple of $p$ (our $a \cdot p$). The second term is our accumulated noise ($2b+m$). Its parity is $m$, and its magnitude is bounded.
    \end{itemize}
\end{frame}

\begin{frame}{Homomorphic Evaluation: Addition}
    Let the Evaluate algorithm simply be integer addition and multiplication. 
    Let $c_1 = q_1 p + 2r_1 + m_1$ and $c_2 = q_2 p + 2r_2 + m_2$.
    \pause

    \vspace{0.3cm}
    \textbf{Homomorphic Addition:}
    $$c_1 + c_2 = p(q_1 + q_2) + 2(r_1 + r_2) + (m_1 + m_2)$$
    \pause
    
    \vspace{0.3cm}
    \textbf{Observation:}
    The structural format $a\cdot p + \text{noise}$ is perfectly preserved! We have a new multiple of $p$, plus $2 \times (\text{new noise})$, plus the sum of the messages. 
    
    \vspace{0.2cm}
    \textbf{Crucially:} The noise grows \textbf{linearly} ($r_1 + r_2$).
\end{frame}

\begin{frame}{Homomorphic Evaluation: Multiplication}
    Now, let us multiply the ciphertexts:
    $$c_1 \cdot c_2 = (q_1 p + 2r_1 + m_1)(q_2 p + 2r_2 + m_2)$$
    \pause
    
    Expanding this algebraically yields:
    $$c_1 \cdot c_2 = p \big( p \cdot q_1 q_2 + q_1(2r_2 + m_2) + q_2(2r_1 + m_1) \big)$$
    $$+ \ 2 \big( 2r_1 r_2 + r_1 m_2 + r_2 m_1 \big) + m_1 m_2$$
    \pause
    
    \vspace{0.3cm}
    \textbf{Observation:}
    The structure is still preserved. We have a massive multiple of $p$ (first term), plus $2 \times (\text{mixed noise})$, plus the product of the messages $m_1 m_2$.
    
    \vspace{0.2cm}
    \textbf{The Crisis:} Look at the dominant noise term: $\mathbf{2r_1 r_2}$. When we multiply ciphertexts, the noise does not add -- it \textbf{multiplies}.
\end{frame}

\subsection{The Multiplicative Bottleneck}

\begin{frame}{The Decryption Constraint \& The Noise Crisis}
    \textbf{Constraint:} 
    Recall our strict rule for correct decryption. The modulo $p$ operation $(c \pmod p)$ will completely destroy the plaintext unless the absolute value of the total accumulated noise is strictly less than $p/2$.
    \pause
    
    \vspace{0.3cm}
    \begin{block}{The Crisis of the Square}
    Because multiplication squares the magnitude of the noise, the noise bit-length doubles with every sequential multiplication. It grows exponentially with the multiplicative depth of the circuit. We will rapidly hit the $p/2$ boundary.
    \end{block}
\end{frame}

\begin{frame}{Formalizing the Bottleneck}
    Exactly how deep can our circuit be before the noise destroys the data?
    \pause
    
    \begin{lemma}[Permitted Polynomials]
    Let $f(x_1, \dots, x_t)$ be the multivariate polynomial computed by a circuit, and let $d$ be its degree. The scheme can correctly evaluate this polynomial as long as:
    $$d \le \frac{\eta - 4 - \log||\vec{f}||_1}{\rho' + 2}$$
    where $||\vec{f}||_1$ is the $l_1$ norm of the coefficient vector of $f$.
    \end{lemma}
    \pause
    
    \begin{block}{Intuition: The Noise Budget}
    \begin{itemize}
        \item Think of $\eta$ (the bit-size of $p$) as our total \textbf{Noise Budget}.
        \item Think of $\rho'$ (the initial noise size) as the \textbf{cost} of one multiplication.
        \item Therefore, the maximum multiplicative depth (degree $d$) we can handle is roughly our total budget divided by the cost per operation: $\approx \eta / \rho'$.
    \end{itemize}
    \end{block}
\end{frame}

\begin{frame}{Story so far}
    \textbf{Where do we stand?}
    \pause
    
    \begin{enumerate}
        \item We have successfully built an encryption scheme that can evaluate additions and multiplications without the server knowing the underlying data.
        \pause
        \item However, because of the multiplicative noise bottleneck, this scheme is strictly a \textbf{Somewhat Homomorphic Encryption (SHE)} scheme.
        \pause
        \item If we evaluate a circuit deeper than $d \approx \eta / \rho'$, the noise wraps around the modulus $p$, and the data is lost forever.
    \end{enumerate}
    \pause
    
    \vspace{0.5cm}
    \begin{center}
        \textit{To achieve our dream of infinite FHE computation, we must find a way to evaluate our own decryption circuit within this strict noise budget. This leads us to the next part - the reduction of the hardness of the problem to a known hard problem, and the squashing of the decryption circuit.}
    \end{center}
\end{frame}

% \section{Variations and Optimizations}

% \subsection{The Practicality Gap}
% \begin{frame}{The Practicality Gap}
%     \begin{motivation}
%     We have established the mathematical engine for our Somewhat Homomorphic Encryption (SHE) scheme. However, before we can secure it or bootstrap it, we face two massive engineering realities:
%     \begin{enumerate}
%         \item \textbf{Memory Explosion:} Ciphertexts grow exponentially in physical bit-size during evaluation.
%         \item \textbf{Bandwidth Choke:} Even fresh ciphertexts are impractically large ($\tilde{\mathcal{O}}(\lambda^5)$ bits).
%     \end{enumerate}
%     \end{motivation}
%     \pause
    
%     \vspace{0.3cm}
%     To make this scheme practically viable, we must introduce two structural optimizations:
%     \begin{itemize}
%         \item \textit{In-flight:} Modular reduction during evaluation.
%         \item \textit{Post-flight:} Ciphertext compression via group exponentiation.
%     \end{itemize}
% \end{frame}

% \subsection{Optimization 1: Modular Reduction During Evaluate}
% \begin{frame}{The Size Explosion Problem}
%     Recall our encryption algorithm: $c = \left[ m + 2r + 2\sum_{i \in S} x_i \right]_{x_0}$.
%     \pause
    
%     \vspace{0.3cm}
%     \textbf{The Problem:}
%     During $\text{Encrypt}$, we easily keep the ciphertext size bounded to $\gamma$ bits by reducing modulo $x_0$. 
    
%     However, during $\text{Evaluate}$, we \textbf{cannot} simply reduce modulo $x_0$ after a multiplication. 
%     \pause
    
%     \vspace{0.3cm}
%     \textbf{Why?} 
%     After one multiplication, the ciphertext squares in size, becoming vastly larger than $x_0$. A standard modulo $x_0$ reduction would subtract a massive multiple of $x_0$. Because $x_0$ contains its own noise ($r_0$), subtracting a massive multiple of $x_0$ injects a catastrophic amount of noise into the ciphertext, instantly destroying the plaintext.
% \end{frame}

% \begin{frame}{The Solution: A Ladder of Zeroes}
%     \textbf{The Fix:} We augment the public key with a "ladder" of increasingly massive encryptions of zero.
%     \pause
    
%     \vspace{0.2cm}
%     For $i = 0, \dots, \gamma$, we generate new elements $x'_i$:
%     $$q'_i \leftarrow \mathbb{Z} \cap [2^{\gamma+i-1}/p, 2^{\gamma+i}/p)$$
%     $$x'_i = 2(q'_i \cdot p + r'_i)$$
%     where $r'_i$ is drawn from the standard noise distribution $(-2^\rho, 2^\rho)$.
%     \pause
    
%     \vspace{0.3cm}
%     \begin{intuition}[The Step-Down Ladder]
%     These $x'_i$ values are astronomically large, scaling up to $2^{2\gamma}$. 
    
%     During evaluation, when a ciphertext exceeds $2^\gamma$, we do not reduce by $x_0$. Instead, we reduce modulo the largest $x'_\gamma$, then modulo $x'_{\gamma-1}$, and so on, stepping down the ladder until the ciphertext is back to $\gamma$ bits.
%     \end{intuition}
%     \pause
    
%     \textbf{Result:} Because we only ever subtract small multiples of these $x'_i$ elements at each step, the accumulated error remains strictly bounded!
% \end{frame}

% \subsection{Optimization 2: Ciphertext Compression}
% \begin{frame}{The Bandwidth Problem}
%     Even with the ladder optimization keeping ciphertexts at $\gamma$ bits, $\gamma$ is still massively large -- specifically $\tilde{\mathcal{O}}(\lambda^5)$ bits. Transmitting this over a network is brutally inefficient.
%     \pause
    
%     \vspace{0.3cm}
%     \textbf{The Goal:} Compress the final output ciphertext to the size of a standard RSA modulus \textit{after} the cloud server has finished all homomorphic evaluations.
% \end{frame}

% \begin{frame}{Compression via Group Exponentiation}
%     \textbf{The Mechanism:}
%     We supplement the public key with the description of a cyclic group $G = \langle g \rangle$ whose order is a multiple of the secret key $p$.
%     \pause
    
%     \vspace{0.3cm}
%     \textbf{The Compression Step:}
%     Instead of sending the massive integer $c$ back to the user, the server computes and transmits:
%     $$c^* = g^c \pmod N$$
%     \pause
    
%     \vspace{0.3cm}
%     \textbf{The Decryption Step:}
%     Because $c^*$ is an element of the group, the user decrypts by computing the Discrete Logarithm (DL):
%     $$m = \left( DL_g(c^*) \pmod p \right) \pmod 2$$
% \end{frame}

% \begin{frame}{Trade-offs of Compression}
%     How is computing the Discrete Logarithm efficient for the user?
%     \pause
    
%     \vspace{0.3cm}
%     \textbf{The Secret Key Modification:}
%     We must generate the secret key $p$ as a \textbf{smooth} number (a product of many small distinct primes). This allows the user to efficiently compute the discrete log using the Pohlig-Hellman algorithm or baby-step giant-step method.
%     \pause
    
%     \vspace{0.3cm}
%     \begin{block}{The Ultimate Trade-off}
%     \begin{itemize}
%         \item \textbf{Pros:} The ciphertext is compressed from $\tilde{\mathcal{O}}(\lambda^5)$ bits down to the size of an RSA modulus ($\sim 2048$ bits).
%         \item \textbf{Cons:} The compressed ciphertext $c^*$ resides in the group $G$. Because we do not have an FHE scheme over this group, \textbf{we completely lose all homomorphic properties.} 
%     \end{itemize}
%     \end{block}
%     \pause
    
%     \vspace{0.2cm}
%     \textit{Conclusion:} This is strictly a post-flight optimization, applied only when computation is complete and data is ready for transit.
% \end{frame}

\section{Security Reduction and Known Attacks}

\subsection{The Approximate-GCD Problem}

\begin{frame}{The Need for Provable Security}
    \begin{motivation}
    We have established a working mathematical engine, but in modern cryptography, we do not simply \textit{guess} that a scheme is secure. We must show that breaking the semantic security of our encryption is mathematically equivalent to solving a known, well-studied hard problem.
    \end{motivation}
    \pause
    
    \vspace{0.3cm}
    For RSA, that problem is integer factorization. For Gentry's 2009 FHE, it is finding short vectors in ideal lattices.
    \pause
    
    \vspace{0.3cm}
    For our integer-based scheme, the foundation is the \textbf{Approximate-GCD Problem}.
\end{frame}

\begin{frame}{Defining Approximate-GCD}
    Recall our noise distribution $\mathcal{D}_{\gamma,\rho}(p)$, which outputs integers of the form $x = pq + r$, where $x$ is roughly $\gamma$ bits long, and the noise $r$ is bounded by $2^\rho$.
    \pause

    \vspace{0.3cm}
    \begin{definition}[Approximate Integer Greatest Common Divisor]
    The $(\rho, \eta, \gamma)$-approximate-gcd problem is: given polynomially many samples $x_0, x_1, \dots, x_\tau$ drawn from $\mathcal{D}_{\gamma,\rho}(p)$ for a randomly chosen $\eta$-bit odd integer $p$, output the hidden integer $p$.
    \end{definition}
    \pause
    
    \vspace{0.3cm}
    \begin{intuition}
    If the noise $r=0$, the problem is trivial. You are given exact multiples of $p$, and a single run of the 2,000-year-old Euclidean algorithm computes $\gcd(x_1, x_2) = p$ instantly. The addition of the noise term $r$ is the sole cryptographic barrier.
    \end{intuition}
\end{frame}

% \section{The Security Reduction (High Level)}

% \begin{frame}{The Search-to-Decision Reduction}
%     The authors provide a formal reduction proving that if you can break the scheme, you can solve Approximate-GCD.
%     \pause
    
%     \vspace{0.2cm}
%     \textbf{The Blueprint of the Reduction:}
%     \begin{enumerate}
%         \item \textbf{The Oracle:} Suppose an attacker $\mathcal{A}$ has a non-negligible advantage in predicting the plaintext bit of a ciphertext.
%         \item \textbf{Extracting the Parity:} The authors prove that $\mathcal{A}$ can be mathematically manipulated into acting as a reliable "Oracle" that outputs the Least Significant Bit (LSB) of the quotient $q_p(z)$ for any noisy integer $z$. (The Leftover Hash Lemma ensures the simulated ciphertexts fed to $\mathcal{A}$ are statistically indistinguishable from real ones).
%         \pause
%         \item \textbf{Binary GCD:} We plug this Oracle into a modified Binary-GCD algorithm. By iteratively asking the Oracle for parity bits, we can systematically peel away the noise terms from our samples without knowing $p$.
%         \item \textbf{Extraction:} After $\mathcal{O}(\gamma)$ iterations, the noise is stripped, and the exact secret key $p$ is revealed.
%     \end{enumerate}
% \end{frame}

\subsection{The Security Reduction (Outline)}

\begin{frame}{The Search-to-Decision Reduction}
    \begin{motivation}
    We must mathematically prove that breaking our scheme is equivalent to solving the Approximate-GCD problem. We do this via a ``Search-to-Decision" reduction.
    \end{motivation}
    \pause
    
    \vspace{0.3cm}
    \begin{theorem}
    Any attack $\mathcal{A}$ with advantage $\epsilon$ on the semantic security of the encryption scheme can be converted into an algorithm $\mathcal{B}$ for solving the $(\rho, \eta, \gamma)$-approximate-gcd problem with success probability at least $\epsilon/2$.
    \end{theorem}
    \pause
    
    \vspace{0.3cm}
    \textbf{The Setup:} 
    Algorithm $\mathcal{B}$ is given random noisy multiples of a hidden integer $p$. Its goal is to find $p$. To do this, $\mathcal{B}$ uses the attacker $\mathcal{A}$ as a tool, feeding $\mathcal{A}$ simulated public keys and ciphertexts to extract information about $p$.
\end{frame}

\begin{frame}{The Mechanism: Oracle and Binary GCD}
    How does $\mathcal{B}$ use the attacker $\mathcal{A}$ to find the secret key $p$?
    \pause

    \vspace{0.3cm}
    \textbf{Step 1: The Learn-LSB Oracle}
    \begin{itemize}
        \item $\mathcal{B}$ constructs randomized, simulated ciphertexts from a noisy integer $z$.
        \item By observing whether $\mathcal{A}$ correctly guesses the encrypted bits, $\mathcal{B}$ can deduce the \textbf{parity of the quotient} $q_p(z)$, without ever knowing $p$!
    \end{itemize}
    \pause

    \vspace{0.3cm}
    \textbf{Step 2: The Binary GCD Extraction}
    \begin{itemize}
        \item $\mathcal{B}$ takes two noisy multiples of $p$ and runs a standard Binary-GCD algorithm on them.
        \item Whenever the algorithm needs to know if a value is even or odd, $\mathcal{B}$ queries the $\text{Learn-LSB}$ Oracle.
        \item \textbf{Result:} Iterating this process strips away the noise terms completely, reducing the values down to the exact secret key $p$.
    \end{itemize}
\end{frame}

\begin{frame}{Guaranteeing Success: The Distribution Match}
    For this reduction to hold, we must guarantee that attacker $\mathcal{A}$ cannot tell that it is being manipulated by $\mathcal{B}$. The simulated ciphertexts must look perfectly real.
    \pause
    
    \vspace{0.3cm}
    \textbf{Why $\mathcal{A}$ is successfully fooled:}
    \begin{enumerate}
        \item \textbf{Quotient Uniformity:} By applying the \textbf{Leftover Hash Lemma}, the sum of the random subsets used to build the simulated ciphertext makes the underlying quotient statistically uniform.
        \pause
        \item \textbf{Noise Drowning:} $\mathcal{B}$ injects a massive random noise term (drawn from $2^{\rho'}$) into the simulation. This completely ``drowns out" and statistically hides any underlying discrepancies in the base noise.
    \end{enumerate}
    \pause
    
    \vspace{0.3cm}
    \begin{intuition}
    Because the simulated data is statistically indistinguishable from real FHE traffic, $\mathcal{A}$'s advantage holds. A rigorous counting argument proves that $\mathcal{B}$ will successfully recover the Approximate-GCD secret key $p$ with probability $\ge \epsilon/2$.
    \end{intuition}
\end{frame}

\subsection{Why Standard Attacks Fail}

\begin{frame}{Algebraic Failures: Continued Fractions}
    If we cannot break the semantic security, can we attack the Approximate-GCD directly?
    
    \vspace{0.2cm}
    \textbf{Attempt 1: Continued Fractions}
    Since $x_1 \approx q_1 p$ and $x_2 \approx q_2 p$, we know the public ratio $x_1 / x_2$ is very close to the unknown rational $q_1 / q_2$. Can we find $q_1/q_2$ in the continued fraction expansion of $x_1/x_2$?
    \pause
    
    \vspace{0.2cm}
    \textbf{The Mathematical Failure:}
    Let's analyze the distance between the public ratio and the target quotient ratio:
    $$ \left| \frac{x_1}{x_2} - \frac{q_1}{q_2} \right| = \left| \frac{q_2(q_1 p + r_1) - q_1(q_2 p + r_2)}{q_2(q_2 p + r_2)} \right| \approx \left| \frac{q_2 r_1 - q_1 r_2}{p \cdot q_2^2} \right| $$
    \pause
    
    \begin{block}{The Cross-Noise Barrier}
    The distance is dominated by the cross-multiplied noise term $(q_2 r_1 - q_1 r_2) / p$. Because this distance is significantly larger than $1/q_2^2$, the target fraction is hopelessly obscured by the noise.
    \end{block}
\end{frame}

\subsection{The Geometric Threat: Lattice Attacks}

\begin{frame}{The Simultaneous Diophantine Approximation (SDA)}
    Because simple algebraic attacks fail, attackers must turn to the geometry of numbers -- specifically lattice reduction algorithms like LLL.
    \pause
    
    \vspace{0.2cm}
    Given samples $x_0, x_1, \dots, x_t$, we construct a $(t+1)$-dimensional SDA lattice $L$ spanned by the rows of matrix $M$:
    $$ M = \begin{pmatrix} 
    2^\rho & x_1 & x_2 & \dots & x_t \\ 
    0 & -x_0 & 0 & \dots & 0 \\ 
    0 & 0 & -x_0 & \dots & 0 \\ 
    \vdots & \vdots & \vdots & \ddots & \vdots \\ 
    0 & 0 & 0 & \dots & -x_0 
    \end{pmatrix} $$
\end{frame}

\begin{frame}{Finding the Target Vector}
    The attacker's goal is to find the hidden vector of quotients $\vec{q} = \langle q_0, q_1, \dots, q_t \rangle$. Let's multiply $\vec{q}$ by $M$:
    
    $$ \vec{v} = \vec{q} \cdot M = \langle q_0 2^\rho, \ q_0 x_1 - q_1 x_0, \ \dots, \ q_0 x_t - q_t x_0 \rangle $$
    \pause
    
    \vspace{0.3cm}
    \textbf{The Critical Cancellation:} 
    Look at the components $q_0 x_i - q_i x_0$. Since $x_i = q_i p + r_i$, we can expand:
    $$ q_0 x_i - q_i x_0 = q_0(q_i p + r_i) - q_i(q_0 p + r_0) $$
    $$ = q_0 q_i p + q_0 r_i - q_i q_0 p - q_i r_0 $$
    $$ = \mathbf{q_0 r_i - q_i r_0} $$
    \pause
    
    \begin{block}{The Threat}
    The secret key $p$ perfectly cancels out! This leaves $\vec{v}$ as a relatively short target vector in the lattice. If an algorithm like LLL can find $\vec{v}$, the attacker recovers the quotients, computes $p \approx x_0 / q_0$, and the scheme is broken.
    \end{block}
\end{frame}

\begin{frame}{Defeating the Lattice}
    Will LLL actually find this target vector $\vec{v}$? This is where our parameter selection defends the scheme.
    \pause
    
    \vspace{0.3cm}
    \textbf{Scenario 1: Attacker uses a small number of samples $t$}
    \begin{itemize}
        \item By Minkowski's bound, a low-dimensional lattice contains exponentially many \textit{spurious} vectors that are much shorter than our target $\vec{v}$. 
        \item The target vector is completely buried in geometric noise.
    \end{itemize}
    \pause
    
    \vspace{0.3cm}
    \textbf{Scenario 2: Attacker uses a massive number of samples $t$}
    \begin{itemize}
        \item The attacker forces $t$ to be huge so $\vec{v}$ is the absolute shortest vector.
        \item However, known lattice reduction algorithms take time $\approx 2^{t/k}$ to find a $2^k$ approximation.
        \item Isolating $\vec{v}$ strictly requires time roughly $\mathbf{2^{\gamma / \eta^2}}$.
    \end{itemize}
\end{frame}

\begin{frame}{Summary of Secure Parameters}
    By setting the bit-length of the public key elements ($\gamma$) to be massively large, we force the lattice dimension to be computationally impossible to reduce. 
    \pause
    
    \vspace{0.3cm}
    \textbf{Final Secure Parameters:}
    \begin{itemize}
        \item $\rho = \mathcal{O}(\lambda)$: Prevents brute-force guessing of the noise $r$.
        \item $\eta = \tilde{\mathcal{O}}(\lambda^2)$: Supports the homomorphic noise budget (depth for bootstrapping).
        \item $\gamma = \tilde{\mathcal{O}}(\lambda^5)$: Specifically ensures that $\gamma = \omega(\eta^2 \log \lambda)$. This makes the time $2^{\gamma / \eta^2}$ super-polynomial, thwarting all SDA and lattice reduction algorithms.
    \end{itemize}
    \pause
    
    \vspace{0.3cm}
    \textit{Conclusion:} While this results in massive public keys, it mathematically guarantees rigorous security based solely on the hardness of the Approximate-GCD over the integers.
\end{frame}

\section{Squashing the Decryption Circuit and Bootstrapping}

% ─────────────────────────────────────────────────────────────────────────────
\subsection{The Bootstrapping Blueprint}
% ─────────────────────────────────────────────────────────────────────────────

\begin{frame}{Going from SHE to FHE: The Bootstrapping Blueprint}
    \textbf{The Problem:}
    Our Somewhat Homomorphic Encryption (SHE) scheme can add and multiply.
    However, every multiplication squares the noise.
    \pause

    \vspace{0.3cm}
    Eventually the noise exceeds the decryption boundary $p/2$, and the
    ciphertext is permanently destroyed.
    How can we evaluate a circuit of arbitrary depth?
    \pause

    \vspace{0.3cm}
    \textbf{The Solution: Bootstrapping (Gentry, 2009)}

    If a scheme can homomorphically evaluate its own decryption circuit
    plus one additional gate, it can evaluate \emph{any} circuit of
    \emph{any} depth.
\end{frame}

\begin{frame}{Defining Bootstrappability}
    \begin{definition}[Augmented Decryption Circuits $\mathcal{D}_E(\lambda)$]
    For a given security parameter $\lambda$, $\mathcal{D}_E(\lambda)$
    consists of two circuits, both taking a secret key and two ciphertexts
    as input:    
    \begin{itemize}
        \item One circuit decrypts both ciphertexts and \textbf{adds} the
              plaintext bits mod~2.
        \item One circuit decrypts both ciphertexts and \textbf{multiplies}
              the plaintext bits mod~2.
    \end{itemize}
    \end{definition}

    \vspace{0.3cm}
    \begin{definition}[Bootstrappable Encryption]
    A homomorphic scheme $E$ is \textbf{bootstrappable} if, for every $\lambda$:
    \[
        \mathcal{D}_E(\lambda) \;\subseteq\; \mathcal{C}_E(\lambda)
    \]
    where $\mathcal{C}_E(\lambda)$ is the set of circuits $E$ can correctly
    evaluate.
    \end{definition}
\end{frame}

\begin{frame}{From Bootstrappability to Full Homomorphism}
    \begin{theorem}[Gentry 2009]
    There is an efficient transformation that, given a bootstrappable scheme
    $E$ and a parameter $d = d(\lambda)$, outputs a scheme $E^{(d)}$ such
    that:
    \begin{enumerate}
        \item $E^{(d)}$ is compact, and
        \item $E^{(d)}$ is correct for all circuits of depth up to $d$.
    \end{enumerate}
    Moreover, $E^{(d)}$ is semantically secure if $E$ is.
    \end{theorem}

    \vspace{0.4cm}
    % \textbf{Our goal for the rest of this talk:}

    Now we prove that our squashed scheme satisfies
    \[
        \mathcal{D}_E \;\subset\; \mathcal{C}(\mathcal{P}_E)
    \]
    so that the above theorem applies and we obtain a fully homomorphic scheme.
\end{frame}

% ─────────────────────────────────────────────────────────────────────────────
\subsection{The Depth Problem}
% ─────────────────────────────────────────────────────────────────────────────

\begin{frame}{The Depth Problem: Why Naive Decryption Fails}
    \textbf{Standard decryption:}
    \[
        m' \;=\; (c \bmod p) \bmod 2
        \;=\; \bigl(c - p\lfloor c/p \rceil\bigr) \bmod 2
    \]
    \pause

    \vspace{0.2cm}
    \textbf{The roadblock:} Computing $\lfloor c/p \rceil$ requires integer
    division. As a boolean circuit the carry chain is roughly $\lambda$
    AND-gate layers deep. After $d$ AND layers, noise $\approx r^{2^d}$.
    With $d = \lambda$ this is catastrophic.
    \pause

    \vspace{0.2cm}
    \textbf{The permitted degree bound:}

    A circuit $C$ belongs to $\mathcal{C}(\mathcal{P}_E)$ if its
    polynomial $f$ satisfies:
    \[
        d \;\leq\; \frac{\eta - 4 - \log \|\hat{f}\|_1}{\rho' + 2}
    \]

    \vspace{0.1cm}
    The division circuit's degree far exceeds this bound.
    We must \textbf{squash} the decryption circuit to bring it within reach.
\end{frame}

% ─────────────────────────────────────────────────────────────────────────────
\subsection{Squashing the Decryption Circuit}
% ─────────────────────────────────────────────────────────────────────────────

\begin{frame}{Squashing: The Strategy}
    \textbf{The key observation:}

    The division $\lfloor c/p \rceil$ is expensive because $p$ is secret.
    If the key owner pre-computes a hint about $1/p$ and publishes it,
    the server can approximate the division using only public information  -- 
    and the remaining homomorphic work becomes shallow.
    \pause

    \vspace{0.3cm}
    \textbf{Three new parameters:}
    \begin{itemize}
        \item $\kappa = \gamma\eta/\rho'$:\quad
              precision bits for the fixed-point approximation of $1/p$
        \item $\theta = \lambda$:\quad
              Hamming weight of the sparse secret selector vector $\vec{s}$
        \item $\Theta = \omega(\kappa \cdot \log\lambda)$:\quad
              total size of the public hint vector $\vec{y}$
    \end{itemize}
    \pause

    \vspace{0.3cm}
    \textbf{The shift:} Move the expensive $1/p$ computation from
    decrypt-time (constrained by the noise budget) to key-generation time
    (unconstrained, done once offline by the key owner).
\end{frame}

\begin{frame}{Squashing: Modifying Key Generation}
    Starting from $\mathsf{sk}^* = p$ and $\mathsf{pk}^*$ of the base scheme:

    \vspace{0.15cm}
    \textbf{Step 1  --  Scaled approximation of $1/p$:}
    \[
        x_p \;\leftarrow\; \lfloor 2^\kappa / p \rceil
    \]
    Computed by the key owner in the clear. After this, $p$ is never needed again.
    \pause

    \vspace{0.15cm}
    \textbf{Step 2  --  Sparse secret selector $\vec{s}$:}

    Choose a random $\Theta$-bit binary vector $\vec{s}$ with Hamming weight
    exactly $\theta$. Set $S = \{i : s_i = 1\}$.
    $\vec{s}$ becomes the \textbf{new secret key}, replacing $p$.
    \pause

    \vspace{0.15cm}
    \textbf{Step 3  --  Public hint vector $\vec{y}$:}

    Choose random $u_i \in [0, 2^{\kappa+1})$ with
    $\sum_{i \in S} u_i \equiv x_p \pmod{2^{\kappa+1}}$.
    Set $y_i = u_i / 2^\kappa$. Publish $\vec{y} = \{y_1, \ldots, y_\Theta\}$.
    Each $y_i \in [0,2)$ with $\kappa$ bits of precision, satisfying:
    \[
        \Bigl[\sum_{i \in S} y_i\Bigr]_2 \;=\; \frac{1}{p} - \Delta_p,
        \qquad |\Delta_p| < 2^{-\kappa}
    \]

    Output: $\mathsf{sk} = \vec{s}$,\quad
    $\mathsf{pk} = (\mathsf{pk}^*,\; y_1, \ldots, y_\Theta)$.
\end{frame}

\begin{frame}{Squashing: Modified Encrypt and Decrypt}
    \textbf{Modified Encrypt:}

    Generate $c^*$ as before. Then for each $i \in \{1,\ldots,\Theta\}$:
    \[
        z_i \;\leftarrow\; [c^* \cdot y_i]_2
    \]
    keeping only $n = \lceil\log\theta\rceil + 3$ bits of precision.
    Output $(c^*,\; z_1, \ldots, z_\Theta)$.
    \pause

    \vspace{0.2cm}
    \textbf{Why the $z_i$ encode the division:}
    \[
        \sum_{i \in S} z_i
        \;\approx\; c^* \sum_{i \in S} y_i
        \;\approx\; \frac{c^*}{p} \pmod{2}
    \]
    So $\lfloor \sum_i s_i z_i \rceil$ approximates $\lfloor c^*/p \rceil$
    without the server ever knowing $p$ or $\vec{s}$.
    \pause

    \vspace{0.2cm}
    \textbf{Modified Decrypt:}
    \[
        m' \;\leftarrow\;
        \Bigl[\, c^* - \Bigl\lfloor \sum_{i=1}^\Theta s_i z_i \Bigr\rceil \,\Bigr]_2
    \]
    A sparse dot product of $\theta$ nonzero terms  --  far shallower than
    integer division.
\end{frame}

% ─────────────────────────────────────────────────────────────────────────────
\subsection{The Complete Encryption Scheme}
% ─────────────────────────────────────────────────────────────────────────────

\begin{frame}{The Complete FHE Scheme: KeyGen}
    \textsc{KeyGen}($\lambda$):
    \begin{enumerate}
        \item $p \;\leftarrow$ random odd $\eta$-bit integer
              \hfill{\small(secret; used here only)}
        \item $\mathsf{pk}^* \;\leftarrow \{x_0,\ldots,x_\tau\}$
              where $x_i = q_i p + 2r_i$
              \hfill{\small(encryptions of zero)}
        \item $x_p \;\leftarrow \lfloor 2^\kappa/p \rceil$
        \item $\vec{s} \;\leftarrow$ random $\Theta$-bit binary vector
              with $\mathrm{wt}(\vec{s}) = \theta$;\quad $S = \{i : s_i = 1\}$
        \item Choose $u_i$ at random subject to
              $\displaystyle\sum_{i\in S} u_i \equiv x_p \pmod{2^{\kappa+1}}$;\quad
              set $y_i = u_i/2^\kappa$
        \item Output $\mathsf{sk} = \vec{s}$ and
              $\mathsf{pk} = \bigl(\mathsf{pk}^*,\; y_1, \ldots, y_\Theta\bigr)$
    \end{enumerate}
\end{frame}

\begin{frame}{The Complete FHE Scheme: Encrypt and Evaluate}
    \textbf{\textsc{Encrypt}($\mathsf{pk}$,\; $m \in \{0,1\}$):}
    \begin{enumerate}
        \item Choose random $T \subseteq \{1,\ldots,\tau\}$
              and $r \in (-2^{\rho'}, 2^{\rho'})$
        \item $c^* \;\leftarrow \bigl(m + 2r + 2\sum_{i\in T} x_i\bigr)
              \bmod x_0$
        \item For each $i$:\; $z_i \;\leftarrow [c^* \cdot y_i]_2$,
              keeping $n = \lceil\log\theta\rceil + 3$ bits of precision
        \item Output $(c^*,\; z_1, \ldots, z_\Theta)$
    \end{enumerate}

    \vspace{0.4cm}
    \textbf{\textsc{Evaluate}($\mathsf{pk}$,\; $C$,\; $c_1,\ldots,c_t$):}
    \begin{enumerate}
        \item Apply the integer addition and multiplication gates of $C$
              to the $c^*$-parts of each ciphertext
        \item Noise grows additively under addition, quadratically under
              multiplication
        \item Recompute $z_i \;\leftarrow [c^*_\mathsf{out} \cdot y_i]_2$
              for the output ciphertext
    \end{enumerate}
\end{frame}

\begin{frame}{The Complete FHE Scheme: Decrypt}
    \textbf{\textsc{Decrypt}($\mathsf{sk}$,\; $(c^*, \vec{z})$):}
    \begin{enumerate}
        \item Sparse dot product:
              $\sigma \;\leftarrow \displaystyle\sum_{i=1}^\Theta s_i \cdot z_i$
              \hfill{\small(only $\theta$ nonzero terms)}
        \item Round: $v \;\leftarrow \lfloor \sigma \rceil$
        \item Output $m' = (c^* - v) \bmod 2$
    \end{enumerate}

    \vspace{0.4cm}
    \textbf{Why decryption is correct:}

    $\sigma$ is within $\tfrac{1}{4}$ of the integer $\lfloor c^*/p \rceil$,
    so rounding always recovers the correct value. The bound $\tfrac{1}{4}$
    comes from two error terms, each at most $\tfrac{1}{16}$:
    \begin{itemize}
        \item Truncation error from $n$-bit precision of $z_i$:
              $|\sum_i s_i \Delta_i| \leq \tfrac{1}{16}$
        \item Approximation error from $x_p \approx 2^\kappa/p$:
              $|c^* \cdot \Delta_p| < \tfrac{1}{16}$
    \end{itemize}

    \vspace{0.2cm}
    Total error $< \tfrac{1}{8} < \tfrac{1}{2}$, so rounding is unambiguous.
\end{frame}

% ─────────────────────────────────────────────────────────────────────────────
\subsection{Proving Bootstrappability}
% ─────────────────────────────────────────────────────────────────────────────

\begin{frame}{The Scheme is Bootstrappable}
    \begin{theorem}
    Let $E$ be the squashed scheme and $\mathcal{D}_E$ its augmented
    decryption circuits. Then $\mathcal{D}_E \subset \mathcal{C}(\mathcal{P}_E)$.
    That is, $E$ is bootstrappable, and by [Gentry09, Theorem 4], $E$ is fully
    homomorphic.
    \end{theorem}
    % \pause

    % \vspace{0.3cm}
    % \textbf{Proof strategy.}
    % Express the decryption equation
    % $m' = (c^* - \lfloor\sum_i s_i z_i\rceil) \bmod 2$
    % as a permitted polynomial in the secret key bits $s_i$, and show its
    % degree is within the bound of Equation~(1).

    % \vspace{0.3cm}
    % The computation decomposes into three steps, each contributing a factor
    % to the total degree:

    % \vspace{0.2cm}
    % \begin{center}
    % \renewcommand{\arraystretch}{1.3}
    % \begin{tabular}{clr}
    %     \hline
    %     \textbf{Step} & \textbf{Computation} & \textbf{Degree in $s_i$} \\
    %     \hline
    %     1 & Selection:\; $a_i = s_i \cdot z_i$ \quad($z_i$ is a plain constant) & $2$ \\
    %     2 & Column Hamming weights via elem.\ sym.\ polynomials & $\leq\theta$ \\
    %     3 & Three-for-two on $n{+}1$ numbers; final rounding & $\leq 32\log^2\!\theta$ \\
    %     \hline
    %     & Total (product of steps) & $64\theta\log^2\!\theta$ \\
    %     \hline
    % \end{tabular}
    % \end{center}
\end{frame}

% \begin{frame}{Completing the Proof}
%     From the previous slide: total degree of the decryption polynomial
%     $\leq 64\theta\log^2\!\theta$.

%     \vspace{0.3cm}
%     With $\theta = \lambda$: total degree $\leq 64\lambda\log^2\!\lambda$.

%     \vspace{0.1cm}
%     Augmented decryption (one extra AND or XOR gate, cf.\ Definition~2.5):
%     degree $\leq 128\lambda\log^2\!\lambda$.
%     \pause

%     \vspace{0.3cm}
%     \textbf{Fitting inside $\mathcal{C}(\mathcal{P}_E)$:}

%     By Remark~3.6, with $\alpha = 128$ and $\beta = 7$ (since
%     $\Theta, \tau < \lambda^7$), it is sufficient to set:
%     \[
%         \eta \;=\; \rho \cdot \Theta(\lambda \log^2\!\lambda)
%     \]
%     This makes the permitted degree bound wide enough to contain the
%     augmented decryption circuit.

%     \vspace{0.3cm}
%     Therefore $\mathcal{D}_E \subset \mathcal{C}(\mathcal{P}_E)$, and
%     by [Gentry09, Theorem 4], $E$ is fully homomorphic. \hfill$\square$
%     \pause

%     \vspace{0.3cm}
%     It remains to prove Step~2, i.e.\ to show how the column Hamming
%     weights $W_{-j}$ can be computed by polynomials of degree $\leq \theta$.
%     This is Lemma~6.3.
% \end{frame}

% \begin{frame}{The Column Trick: Why It Is Needed}
%     \textbf{Why not apply three-for-two directly to all $\Theta$ numbers $a_i$?}

%     That gives degree $\leq 32\Theta^{1.71} \approx 32\lambda^{12}$  -- 
%     completely unworkable.
%     \pause

%     \vspace{0.3cm}
%     \textbf{The column trick.}
%     Write each $a_i$ in binary: $a_i = \sum_j 2^{-j} a_{i,-j}$.
%     For each bit position $j$, define the \emph{column Hamming weight}:
%     \[
%         W_{-j} \;=\; \sum_{i=1}^\Theta a_{i,-j}
%     \]
%     Since at most $\theta$ of the $a_i$ are nonzero, $W_{-j} \leq \theta$
%     and fits in $\lceil\log(\theta+1)\rceil < n$ bits. The key identity:
%     \[
%         \sum_i a_i \;=\; \sum_j 2^{-j} W_{-j}
%     \]
%     compresses $\Theta$ numbers into $n+1 \approx \log\theta$ numbers.
%     Three-for-two then runs on $\log\theta$ inputs instead of $\Theta$,
%     giving degree $\leq 32\log^2\!\theta$ instead of $32\lambda^{12}$.
% \end{frame}

% \begin{frame}{Hamming Weight Bits as Polynomials}
%     \begin{lemma}
%     Let $\vec{\sigma} = \langle\sigma_1,\ldots,\sigma_t\rangle$ be a binary
%     vector and $W = W(\vec{\sigma})$ its Hamming weight, with binary
%     representation $W = \sum_{k} 2^k W_k$. Then for every $k$:
%     \[
%         W_k(\vec{\sigma})
%         \;=\; e_{2^k}(\vec{\sigma}) \bmod 2
%         \;=\; \Bigl(\!\sum_{|S|=2^k}\prod_{j \in S} \sigma_j\Bigr) \bmod 2
%     \]
%     This has degree exactly $2^k$ in $\sigma_1,\ldots,\sigma_t$.
%     Moreover, all bits $W_0,\ldots,W_k$ can be computed simultaneously by
%     an arithmetic circuit of size $2^k \cdot t$.
%     \end{lemma}
%     \pause

%     \vspace{0.2cm}
%     \textbf{Why degree $\leq \theta$ in our setting:}

%     In Step~2 the inputs $\sigma_j = a_{i,-j}$ have at most $\theta$
%     nonzero entries. So $e_k(\vec{\sigma}) = 0$ for all $k > \theta$,
%     meaning we only ever need $e_1, e_2, \ldots, e_\theta$.
%     The degree is therefore at most $\theta$, regardless of $t = \Theta$.

%     \vspace{0.2cm}
%     \textbf{Note:} The circuit above is not shallow. But
%     since it computes low-degree polynomials, Lemma~3.5 guarantees it is
%     a permitted circuit.
% \end{frame}

% ─────────────────────────────────────────────────────────────────────────────
\subsection{Security of the Squashed Scheme}
% ─────────────────────────────────────────────────────────────────────────────

\begin{frame}{Security: Two Hardness Assumptions}
    The squashed scheme requires two independent hardness assumptions.
    \pause

    \vspace{0.3cm}
    \textbf{Assumption 1  --  Approximate GCD (AGCD):}

    Given many samples $x_i = q_i \cdot p + r_i$ where $r_i$ is small
    and $p$ is a hidden large odd integer, find $p$.

    \vspace{0.1cm}
    Defends the base scheme: public key elements and fresh ciphertexts
    are computationally indistinguishable from random integers.
    \pause

    \vspace{0.3cm}
    \textbf{Assumption 2  --  Sparse Subset Sum (SSSP):}

    Given $y_1, \ldots, y_\Theta$ and the knowledge that some
    $\theta$-sparse subset sums to $\approx 1/p \pmod{2}$,
    find the subset $S$.

    \vspace{0.1cm}
    Defends the hint vector $\vec{y}$: knowing $\vec{y}$ reveals nothing
    about which $\theta$ indices form $S$.
    \pause

    \vspace{0.3cm}
    \textbf{Note:} SSSP is a separate assumption, not implied by AGCD.
\end{frame}

\section{Evolution and Open Problems}

\begin{frame}{The Evolution of Integer FHE}
    The original 2010 DGHV scheme was a conceptual breakthrough, but suffered from massive public key sizes ($\tilde{\mathcal{O}}(\lambda^{10})$) and severe performance bottlenecks. Subsequent research rapidly optimized the integer paradigm:
    \pause

    \vspace{0.2cm}
    \textbf{1. Key Compression \& Modulus Switching (2011--2012):}
    \begin{itemize}
        \item Evaluated quadratic forms and adopted ``modulus switching" (adapted from the BGV scheme) to slash the public key size down to $\tilde{\mathcal{O}}(\lambda^5)$.
        \item \textit{Crucial impact:} Completely eliminated the need for the complex ``squashing" step prior to bootstrapping.
    \end{itemize}
    \pause

    \textbf{2. Batching \& SIMD Operations (2013):}
    \begin{itemize}
        \item Utilized the Chinese Remainder Theorem (CRT) with multiple secret primes to encrypt vectors of plaintext bits into a single ciphertext, drastically improving amortized efficiency.
    \end{itemize}
    \pause

    \textbf{3. Scale-Invariant Noise \& Message Spaces (2014--2015):}
    \begin{itemize}
        \item Achieved linear (rather than exponential) noise growth, dropping the modulus switching requirement entirely.
        \item Expanded the message space from binary ($\mathbb{Z}_2$) to any prime ($\mathbb{Z}_p$), reducing the multiplicative degree of the decryption circuit.
    \end{itemize}
\end{frame}

\begin{frame}{Open Problems and the Future of FHE}
    While the integer-based approach has been heavily optimized, several critical challenges remain for the Fully Homomorphic Encryption community at large:
    \pause

    \vspace{0.3cm}
    \textbf{1. The Bootstrapping Bottleneck:}
    \begin{itemize}
        \item Even with scale-invariant noise and SIMD batching, homomorphically evaluating the decryption circuit remains a massive computational tax. Minimizing or entirely bypassing the need for bootstrapping remains the holy grail of FHE.
    \end{itemize}
    \pause

    \vspace{0.2cm}
    \textbf{2. Ciphertext Expansion Ratio:}
    \begin{itemize}
        \item The ratio of the ciphertext size to the actual underlying plaintext data is still highly impractical for low-bandwidth environments or massive databases.
    \end{itemize}
    \pause

    \vspace{0.2cm}
    \textbf{3. RLWE vs. Approximate-GCD:}
    \begin{itemize}
        \item Modern practical FHE libraries (like HElib, SEAL, and TFHE) have largely standardized on Ring Learning With Errors (RLWE) over polynomial rings. 
        \item \textit{Open Question:} Can the conceptual simplicity of the integer approach (Approximate-GCD) ever be engineered to competitively match the hardware-accelerated speeds of RLWE schemes?
    \end{itemize}
\end{frame}

\begin{frame}{Conclusion}
    \begin{center}
        \Large \textbf{Thank You!}
        
        \vspace{0.5cm}
        \normalsize
        Are there any questions?
    \end{center}
\end{frame}

\end{document}